% !PRINTED 630
\weberpage{630}
%<<CONT
mit $\mathfrak p_1$ conjugirten Primfactoren theilbar, und es folgt, da $p$ nicht
durch $\mathfrak p^2$ theilbar ist:

1. \gesperrt{Der Primfactor $\mathfrak p_1$ ist der grösste gemeinschaftliche
Theiler von $p$ und $P_1(r)$; in gleicher Weise ergiebt sich, dass
$\mathfrak p_n$ der grösste gemeinschaftliche Theiler von $p$ und $P_1(r^n)$
ist, oder auch der grösste gemeinschaftliche Theiler von $p$ und
$P_{n'}(r)$, wenn $n n' \equiv 1 \ (\mathrm{mod}\ m)$ ist.}

Das letztere ergiebt sich aus der aus (4) durch die Substitution
 $(r, r^{n'})$ folgenden Congruenz
\[
F_1(t) \equiv P_n(t) \ (\mathrm{mod}\ \mathfrak p_{n'}).
\]

Man erhält also die sämmtlichen Primfactoren von $p$, wenn
man die grössten gemeinschaftlichen Theiler von $p$ mit jeder
der rationalen Functionen
\weq{10}{P_{\xi_1}(r),\ P_{\xi_2}(r),\ \ldots,\ P_{\xi_e}(r)}
aufsucht, und zwar ist $\mathfrak p_{\xi_i}$ der grösste gemeinschaftliche Theiler
von $p$ und $P_{\xi_i^{-1}}$.
% [sic] Source omits the argument (r) after P_{xi_i^{-1}} here; (10) and theorem 1 give the intended evaluation.

Wir wollen den Fall noch etwas näher betrachten, wo $f=1$
ist, also
\weq{11}{p \equiv 1 \ (\mathrm{mod}\ m).}

In diesem Falle ist $e = \varphi(m)$; die Gruppe $\mathfrak A$ reducirt sich
auf die Einheit, und die sämmtlichen $\mu$ conjugirten Factoren $\mathfrak p$
von $p$ sind von einander verschieden. Die Formen
\[
P_{\xi_1}(t),\ P_{\xi_2}(t),\ \ldots,\ P_{\xi_e}(t)
\]
werden alle linear, d.\,h. $r$ ist einer rationalen Zahl nach jedem
der Moduln $\mathfrak p_i$ congruent.

Dies ergiebt sich auch daraus, dass hier die Anzahl der
nach dem Modul $\mathfrak p$ incongruenten Zahlen gleich $N(\mathfrak p) = p$ ist, und
dass also $0, 1, 2, \ldots, p-1$ ein volles Restsystem ist.

Ist hiernach etwa $r \equiv c \ (\mathrm{mod}\ \mathfrak p)$, so muss, da $r^m = 1$ ist,
$c^m \equiv 1 \ (\mathrm{mod}\ \mathfrak p)$ also auch $(\mathrm{mod}\ p)$ sein, und wenn also $g$ eine
primitive Wurzel der Primzahl $p$ ist, so muss
\weq{12}{c \equiv g^{-n\frac{p-1}{m}} \ (\mathrm{mod}\ p)}
sein, worin $n$ relativ prim zu $m$ ist. Lassen wir $\mathfrak p$ das System
der conjugirten Ideale durchlaufen, so muss $n$ in (12) die
Gruppe $\mathfrak N$ durchlaufen, und es ist also nur Sache der Bezeichnung,
wenn wir festsetzen:
\weq{13}{r \equiv g^{-\frac{p-1}{m}} \ (\mathrm{mod}\ \mathfrak p).}
